\section{Implementation Considerations}
\label{sec:implementation}

The Semut Adaptive Architecture is intended as a conceptual systems architecture rather than a software framework tied to a particular implementation. Consequently, the proposed memory structures may be implemented using a variety of databases, workflow engines, reasoning systems, and execution environments.

\subsection{Artifact Memory}

Artifact Memory may be implemented using conventional document databases, relational databases, object storage systems, vector databases, or combinations thereof. Depending on the application domain, artifacts may include documents, source code, emails, conversations, configuration files, images, structured records, or other persistent information.

Retrieval mechanisms are similarly implementation dependent. Systems may employ keyword search, semantic retrieval, metadata filtering, graph traversal, or hybrid retrieval strategies. The proposed architecture requires only that relevant contextual information be retrievable before reasoning begins.

\subsection{Action Graph}

The Action Graph represents executable behavior independently from the underlying execution environment. Actions may correspond to API calls, software functions, robotic operations, human tasks, workflow nodes, or external services.

Graph edges represent observed relationships between actions, including ordering constraints, dependencies, alternative execution paths, and historical transitions. As additional executions occur, the graph continuously evolves by incorporating newly observed relationships and execution outcomes.

\subsection{Feedback Memory}

Feedback Memory aggregates observations produced throughout system operation. Feedback may originate from multiple sources, including human approval, software-generated metrics, execution status, application logs, environmental sensors, or explicit user corrections.

Rather than serving exclusively as monitoring data, Feedback Memory becomes part of the adaptive memory available during future reasoning. This enables previously successful behavior, recurring failures, and accumulated user preferences to influence subsequent decision making.

\subsection{Reasoning Engine Independence}

A defining characteristic of the proposed architecture is the separation between adaptive memory and reasoning. The architecture therefore does not prescribe a particular reasoning algorithm.

Possible reasoning engines include:

\begin{itemize}
\item Rule-based systems
\item Large Language Models
\item Classical planning algorithms
\item Reinforcement learning systems
\item Future foundation models
\item Human decision makers
\end{itemize}

Each reasoning engine interacts with the same adaptive memory through retrieval and memory update operations, allowing reasoning algorithms to evolve independently from the underlying memory architecture.

\subsection{Execution Layer}

Execution is likewise independent of the adaptive memory. Actions may be performed through software APIs, workflow platforms, command-line tools, robotic systems, web applications, or direct human interaction.

This separation allows the proposed architecture to coordinate heterogeneous software environments without requiring modifications to existing execution platforms.

Overall, the Semut Adaptive Architecture provides a common adaptive interface connecting persistent knowledge, executable behavior, and accumulated feedback while remaining independent of the specific technologies used to implement each component.
